\section{Introduction}

Sorting waste at the source requires a small decision that is repeated many times: which bin each object goes into. An assistant that sees the object and listens to the person can support that decision, but it needs models from two different domains. Computer vision has pretrained feature extractors that can be adapted with few images. Speech recognition must, in addition, solve the alignment between a long signal and a short text whose temporal correspondence is unknown.

This work presents \emph{EcoVoz}, a local application that integrates three modalities: images, video classified frame by frame, and audio transcribed to text. We address three questions:
\begin{enumerate}
\item how much does each level of fine-tuning contribute in transfer learning, and is a large network (ResNet-50) justified over a mobile one (MobileNetV3-Large)?
\item what can a CTC model learn from only 100 recordings, and how much do data augmentation and pretraining help?
\item is a desktop computer enough to respond in near real time across the three modalities?
\end{enumerate}

The contributions are a comparison of eight vision configurations on a single object-grouped split, a series of five speech experiments evaluated on the same unseen speakers, an end-to-end latency measurement on GPU and CPU, and a repository with the code that produces every figure in this paper.

The rest of the paper is organized as follows: Section~\ref{sec:teoria} summarizes the theoretical background, Section~\ref{sec:metodo} describes the data, models, training, and application, Section~\ref{sec:resultados} analyzes the results, and Section~\ref{sec:conclusiones} concludes.
